\documentclass[11pt,a4paper]{article}

\usepackage[utf8]{inputenc}
\usepackage[T1]{fontenc}
\usepackage{amsmath,amssymb,amsthm}
\usepackage{algorithm}
\usepackage{algpseudocode}
\usepackage{booktabs}
\usepackage{enumitem}
\usepackage{microtype}
\usepackage[margin=2.5cm]{geometry}
\usepackage{hyperref}
\hypersetup{colorlinks=true,linkcolor=blue,citecolor=blue}

% ── Notation ──────────────────────────────────────────────────────────
\newcommand{\G}{\mathcal{G}}
\newcommand{\T}{\mathcal{T}}
\newcommand{\E}{\mathcal{E}}
\newcommand{\R}{\mathcal{R}}
\newcommand{\Types}{\mathrm{Types}}
\newcommand{\range}{\mathrm{range}}
\newcommand{\True}{\mathsf{true}}
\newcommand{\admiss}{\mathrm{admissible}}

\DeclareMathOperator*{\argmax}{arg\,max}

% ── Theorems ──────────────────────────────────────────────────────────
\newtheorem{definition}{Definition}
\newtheorem{proposition}{Proposition}

\title{Graph-Constrained Reasoning with Symbolic Type Gates:\\
The GCR Baseline and DCA-Trie Algorithms}
\author{Bernard Katamanso}
\date{\today}

\begin{document}
\maketitle

%======================================================================
\section{Introduction}
%======================================================================

Large language models (LLMs) can reason over knowledge graphs (KGs) by
generating natural language reasoning paths that traverse edges in the
graph.  \emph{Graph-constrained decoding} (GCD) enforces this traversal
by restricting the LLM's vocabulary at each decode step to tokens that
correspond to valid edges in the KG.

This technical report describes two algorithms for graph-constrained
reasoning:

\begin{enumerate}[label=\textbf{A\arabic*},leftmargin=2.5em]
  \item \textbf{GCR Baseline.}  Single-pass constrained decoding over
    a pre-built trie containing all valid paths up to $L$ hops.
  \item \textbf{DCA-Trie.}  Two variants that add symbolic type gates
    (from a TypeOracle) to constrain which edges enter the trie:
    \begin{enumerate}[label=(\alph*)]
      \item \textbf{v1 Static.}  Pre-filters all paths through the
        TypeOracle before building the trie.
      \item \textbf{v2 Iterative.}  Hop-by-hop trie expansion with
        beam search, applying type gates at each step.
    \end{enumerate}
\end{enumerate}

Both algorithms share the same underlying machinery---a trie that
constrains the LLM's token vocabulary---but differ in \emph{which
paths enter the trie} and \emph{how many forward passes} are required.

%======================================================================
\section{Preliminaries}
%======================================================================

\subsection{Knowledge Graph}

\begin{definition}[Knowledge Graph]
A knowledge graph is a tuple $\G = (V, E, \R)$ where $V$ is a set of
entities, $\R$ is a set of relation types, and $E \subseteq V \times
\R \times V$ is a set of labeled edges.
\end{definition}

Given a natural language question $q$, a KG subgraph $\G_q \subseteq
\G$ is extracted containing entities mentioned in $q$ and their
neighbours up to $L$ hops.  The question entities are $V_q \subseteq
V$.

\subsection{Reasoning Path}

\begin{definition}[Reasoning Path]
A reasoning path of length $L$ through $\G$ is a chain
\[
  c = (v_0, e_1, v_1, e_2, v_2, \ldots, e_L, v_L)
\]
where $v_0 \in V_q$ and $(v_{l-1}, e_l, v_l) \in E$ for all $l
\in \{1, \ldots, L\}$.
\end{definition}

The terminal entity $v_L$ is the candidate answer.

\subsection{Trie}

\begin{definition}[Trie]
A trie $\T$ is a prefix tree over token sequences.  Given a tokenizer
$\phi$, each reasoning path $c$ is encoded as a token sequence
$\phi(c) = (t_1, t_2, \ldots, t_n)$.  The trie stores all such
sequences and supports prefix-constrained generation: at each decode
step with prefix $p$, the set of valid next tokens is
\[
  \T(p) = \{ t \mid \exists\, \text{sequence in } \T \text{ starting
  with } p \cdot t \}.
\]
\end{definition}

\subsection{Constrained Decoding}

Given a prompt $\mathbf{x}$ encoding the question $q$, an LLM produces
token probabilities $P(t_l \mid t_{<l}, \mathbf{x})$ at each step
$l$.  Constrained decoding modifies this to
\[
  P_{\T}(t_l \mid t_{<l}, \mathbf{x}) =
  \begin{cases}
    \frac{P(t_l \mid t_{<l}, \mathbf{x})}{\sum_{t' \in \T(t_{<l})}
    P(t' \mid t_{<l}, \mathbf{x})} & \text{if } t_l \in \T(t_{<l}), \\
    0 & \text{otherwise}.
  \end{cases}
\]

This is implemented via a \texttt{prefix\_allowed\_tokens\_fn} callback
that returns $\T(p)$ at each decode step, masking out all other tokens
with $-\infty$ log-probability.

%======================================================================
\section{Algorithm~1: GCR Baseline}
%======================================================================

The GCR baseline (Luo et al., 2025) performs single-pass constrained
decoding over a trie built from \emph{all} reachable paths.

\subsection{Procedure}

\begin{enumerate}
  \item \textbf{Subgraph extraction.}  Extract $\G_q$ from $\G$ using
    the question entities $V_q$.
  \item \textbf{Path enumeration.}  Run depth-first search (DFS) from
    each $v \in V_q$ up to $L$ hops, collecting all simple paths.
  \item \textbf{Trie construction.}  Tokenize each path string and
    insert into a MarisaTrie.
  \item \textbf{Constrained generation.}  Run the LLM with the trie
    constraint, generating the full reasoning path and answer in one
    forward pass.
\end{enumerate}

\begin{algorithm}[H]
\caption{GCR Baseline: Single-Pass Constrained Decoding}
\label{alg:baseline}
\begin{algorithmic}[1]
\Require question $q$, entities $V_q$, KG subgraph $\G_q$, max hops
  $L$, LLM $\mathcal{M}$, tokenizer $\phi$
\Ensure reasoning path $c^*$, answer $a$
\State $\mathcal{P} \gets \textsc{DFS}(\G_q, V_q, L)$
  \Comment{enumerate all paths}
\State $\T \gets \textsc{BuildTrie}(\phi, \mathcal{P})$
  \Comment{insert all path tokens}
\State $\mathbf{x} \gets \textsc{Prompt}(q)$
  \Comment{encode question as prompt}
\State $c^* \gets \mathcal{M}.\textsc{Generate}(\mathbf{x}, \T)$
  \Comment{single-pass constrained decoding}
\State $a \gets \textsc{ExtractAnswer}(c^*)$
\State \Return $c^*, a$
\end{algorithmic}
\end{algorithm}

\subsection{Properties}

\begin{itemize}
  \item \textbf{Completeness.}  Every reachable path up to $L$ hops is
    in the trie.  The LLM can choose any of them.
  \item \textbf{Single forward pass.}  The entire path is generated in
    one call to $\mathcal{M}.\textsc{Generate}$.
  \item \textbf{No semantic filtering.}  All topologically valid paths
    are admitted, regardless of type compatibility.
  \item \textbf{Trie size.}  $|\mathcal{P}| = O(d^L)$ paths, where $d$
    is the average out-degree.  For typical KG subgraphs ($d \approx
    20$, $L = 2$), this yields thousands of paths.
\end{itemize}

%======================================================================
\section{Algorithm~2: DCA-Trie}
%======================================================================

DCA-Trie adds a \emph{TypeOracle} that applies symbolic type gates to
filter paths before they enter the trie.  The key insight: not all
topologically valid paths are semantically coherent.

\subsection{The TypeOracle}

The TypeOracle is a purely symbolic component that extracts type
information from the KG's ontology schema (Freebase
\texttt{common.topic.notable\_types}, \texttt{rdf-schema\#range},
etc.).  It provides two gates:

\begin{definition}[Range Gate]
Given a relation $e \in \R$ and a tail entity $v$, the range gate
checks whether $v$'s types are compatible with $e$'s declared range:
\[
  \range\textsc{Gate}(e, v) \;\iff\;
  \T(v) \cap \range(e) \neq \emptyset
  \;\lor\; \T(v) = \emptyset
  \;\lor\; \range(e) = \emptyset.
\]
The empty-set disjuncts make the gate \emph{conservative}: if types
are unknown, the path is admitted.
\end{definition}

\begin{definition}[Type Gate]
Given a candidate terminal entity $v$, expected answer types $A
\subseteq \Types$, hop index $l$, and max hops $L$, the type gate
checks:
\[
  \textsc{TypeGate}(v, A, l, L) \;\iff\;
  l < L
  \;\lor\; A = \emptyset
  \;\lor\; \T(v) = \emptyset
  \;\lor\; \T(v) \cap A \neq \emptyset.
\]
At intermediate hops ($l < L$), the gate always passes.  At the
terminal hop, it checks type compatibility with the expected answer
types.
\end{definition}

\begin{definition}[Admissibility]
A candidate edge $(e_l, v_l)$ is \emph{admissible} iff it passes both
gates:
\[
  \admiss(e_l, v_l) \;=\; \range\textsc{Gate}(e_l, v_l)
  \;\land\; \textsc{TypeGate}(v_l, A, l, L).
\]
\end{definition}

\subsection{Answer Type Inference}

The expected answer types $A$ are inferred from the question string via
pattern matching:
\[
  A = \bigcup_{\text{pattern}(q) \to T_k} T_k
\]
where $T_k$ are pre-defined type sets (e.g., ``\emph{who}'' $\to$
$\mathrm{PersonTypes}$, ``\emph{where}'' $\to$
$\mathrm{LocationTypes}$).  If no pattern matches, $A = \emptyset$
(unconstrained), and the type gate passes everything.

\subsection{Algorithm~2a: DCA-Trie v1 (Static)}

v1 is a direct refinement of the GCR baseline: filter paths through
the TypeOracle \emph{before} building the trie.

\begin{algorithm}[H]
\caption{DCA-Trie v1: Static Type-Oracle Filtering}
\label{alg:v1}
\begin{algorithmic}[1]
\Require question $q$, entities $V_q$, KG subgraph $\G_q$, max hops
  $L$, LLM $\mathcal{M}$, tokenizer $\phi$, TypeOracle $\Omega$
\Ensure reasoning path $c^*$, answer $a$
\State $\mathcal{P} \gets \textsc{DFS}(\G_q, V_q, L)$
  \Comment{enumerate all paths}
\State $A \gets \Omega.\textsc{InferAnswerTypes}(q)$
  \If{$A = \emptyset$}
    \State $A \gets \Omega.\textsc{InferFromPaths}(\mathcal{P})$
    \Comment{fallback: use terminal entity types}
  \EndIf
\State $\mathcal{P}_{\text{filt}} \gets \{\}$
\For{$c = (v_0, e_1, v_1, \ldots, e_L, v_L) \in \mathcal{P}$}
  \If{$\forall\, l:\; \range\textsc{Gate}(e_l, v_l)$ \textbf{and}
    $\textsc{TypeGate}(v_L, A, L, L)$}
    \State $\mathcal{P}_{\text{filt}} \gets \mathcal{P}_{\text{filt}}
      \cup \{c\}$
  \EndIf
\EndFor
\State $\T \gets \textsc{BuildTrie}(\phi, \mathcal{P}_{\text{filt}})$
\State $\mathbf{x} \gets \textsc{Prompt}(q)$
\State $c^* \gets \mathcal{M}.\textsc{Generate}(\mathbf{x}, \T)$
\State $a \gets \textsc{ExtractAnswer}(c^*)$
\State \Return $c^*, a$
\end{algorithmic}
\end{algorithm}

\paragraph{Properties.}
\begin{itemize}
  \item \textbf{Same single forward pass} as the baseline.
  \item \textbf{Smaller trie.}  Filtering removes type-incompatible
    paths, typically reducing trie size by 25\% (WebQSP) to 40\% (CWQ).
  \item \textbf{Risk: over-filtering.}  If the TypeOracle is too
    aggressive, correct paths are removed and the answer becomes
    unreachable.  This is measured by the \emph{false negative rate}
    (FNR) of the oracle.
\end{itemize}

\subsection{Algorithm~2b: DCA-Trie v2 (Iterative Beam Search)}

v2 replaces the single-pass approach with hop-by-hop trie expansion.
At each hop, the trie is rebuilt from the 1-hop neighbours of the
current beam's head pool, filtered through the TypeOracle.

\begin{definition}[Head Pool]
For a beam $b$ at step $l$, the head pool $H_b^{(l)} \subseteq V$ is
the set of entities that are available as starting points for the next
hop.  Initially $H_b^{(0)} = V_q$.  After each hop, the new entity
$v_l$ is added: $H_b^{(l)} = H_b^{(l-1)} \cup \{v_l\}$.
\end{definition}

\begin{definition}[Beam Unit]
A beam unit is a tuple $b = (\sigma, H_b, s, l)$ where:
\begin{itemize}
  \item $\sigma$ is the accumulated prompt (question + generated path
    so far),
  \item $H_b$ is the head pool,
  \item $s$ is the cumulative log-probability score,
  \item $l$ is the current hop index.
\end{itemize}
\end{definition}

\begin{algorithm}[H]
\caption{DCA-Trie v2: Iterative Beam Search with Type-Oracle Gates}
\label{alg:v2}
\begin{algorithmic}[1]
\Require question $q$, entities $V_q$, KG subgraph $\G_q$, max hops
  $L$, LLM $\mathcal{M}$, tokenizer $\phi$, TypeOracle $\Omega$,
  beam size $B$, max new tokens $N$
\Ensure reasoning path $c^*$, answer $a$
\State $A \gets \Omega.\textsc{InferAnswerTypes}(q)$
\State $\mathbf{x} \gets \textsc{Prompt}(q)$
\State $\mathcal{B}_0 \gets \{\}$
\For{$v \in V_q$}
  \State $\mathcal{P}_1 \gets
    \textsc{GatedPaths}(\G_q, v, \Omega, A, 1, L)$
  \If{$\mathcal{P}_1 \neq \emptyset$}
    \State $\mathcal{B}_0 \gets \mathcal{B}_0 \cup
      \{(\mathbf{x}, \{v\}, 0, 0)\}$
  \EndIf
\EndFor
\If{$\mathcal{B}_0 = \emptyset$} \Return \textsc{None} \EndIf
\State $\mathcal{B} \gets \mathcal{B}_0$
\For{$l = 1$ \textbf{to} $L$}
  \State $\mathcal{B}' \gets \{\}$
  \For{beam $b = (\sigma, H_b, s, l{-}1) \in \mathcal{B}$}
    \State $\mathcal{P}_l \gets \{\}$
    \For{$h \in H_b$}
      \State $\mathcal{P}_l \gets \mathcal{P}_l \cup
        \textsc{GatedPaths}(\G_q, h, \Omega, A, l, L)$
    \EndFor
    \If{$\mathcal{P}_l = \emptyset$} \textbf{continue} \EndIf
    \State $\T_b \gets \textsc{BuildTrie}(\phi, \mathcal{P}_l)$
    \State $\sigma' \gets \sigma \;\|\; \texttt{<PATH>}$
    \State $\textsc{Results} \gets \mathcal{M}.\textsc{Generate}(
      \sigma', \T_b, \text{num\_beams}=B, \text{return\_scores}=\True)$
    \For{$(\text{path}_i, s_i) \in \textsc{Results}$}
      \State $v_l \gets \textsc{LastEntity}(\text{path}_i)$
      \State $\sigma'' \gets \sigma \;\|\; \text{path}_i$
      \State $\mathcal{B}' \gets \mathcal{B}' \cup
        \{(\sigma'', H_b \cup \{v_l\}, s + s_i, l)\}$
    \EndFor
  \EndFor
  \State $\mathcal{B} \gets \textsc{TopB}(\mathcal{B}', B)$
  \Comment{keep best $B$ beams by score}
\EndFor
\State $b^* \gets \argmax_{b \in \mathcal{B}} b.s$
\State $c^* \gets \textsc{ParsePath}(b^*.{\sigma})$
\State $a \gets \textsc{ExtractAnswer}(c^*)$
\State \Return $c^*, a$
\end{algorithmic}
\end{algorithm}

\paragraph{The GatedPaths subroutine.}

\begin{algorithm}[H]
\caption{GatedPaths: 1-Hop Enumerate + Type-Oracle Filter}
\label{alg:gated}
\begin{algorithmic}[1]
\Require head entity $h$, KG $\G_q$, TypeOracle $\Omega$, answer types
  $A$, hop $l$, max hops $L$
\Ensure list of allowed 1-hop path strings
\State $\mathcal{P} \gets \{\}$
\For{$(h, e, v) \in \G_q.\textsc{Neighbors}(h)$}
  \If{$\range\textsc{Gate}(e, v)$ \textbf{and}
    ($l < L$ \textbf{or} $\textsc{TypeGate}(v, A, l, L)$)}
    \State $\mathcal{P} \gets \mathcal{P} \cup \{h \to e \to v\}$
  \EndIf
\EndFor
\State \Return $\mathcal{P}$
\end{algorithmic}
\end{algorithm}

\paragraph{Properties.}
\begin{itemize}
  \item \textbf{Multiple forward passes.}  Each hop requires a separate
    call to $\mathcal{M}.\textsc{Generate}$.  With $L$ hops and $B$
    beams, worst case is $L \times B$ forward passes.
  \item \textbf{Adaptive trie size.}  The trie is rebuilt at each hop
    from only the 1-hop neighbours of the current head pool, so it is
    much smaller than the baseline's full-path trie.
  \item \textbf{Beam search.}  Multiple candidate paths are explored
    simultaneously, with scores accumulated across hops.
  \item \textbf{Backtracking.}  If a beam's head pool yields no valid
    paths, the beam is dropped (not expanded further).
\end{itemize}

%======================================================================
\section{Complexity Analysis}
%======================================================================

\subsection{Path Explosion}

The number of simple paths of length $L$ from a set of start entities
in a graph with average out-degree $d$ is
\[
  |\mathcal{P}| = O(|V_q| \cdot d^L).
\]
For WebQSP ($|V_q| \approx 1$, $d \approx 20$, $L = 2$): $|\mathcal{P}|
\approx 400$ paths.  For CWQ ($d \approx 30$, $L = 3$):
$|\mathcal{P}| \approx 27{,}000$ paths.

\subsection{Comparison}

\begin{table}[H]
\centering
\caption{Complexity comparison of the three algorithms.}
\label{tab:complexity}
\begin{tabular}{@{}lccc@{}}
\toprule
 & \textbf{GCR Baseline} & \textbf{DCA-Trie v1} &
  \textbf{DCA-Trie v2} \\
\midrule
DFS passes & 1 & 1 & 0 \\
Trie builds & 1 & 1 & $L$ \\
Model forward passes & 1 & 1 & $\leq L \times B$ \\
Trie size (paths) & $O(d^L)$ & $O(d^L)$ & $O(d)$ \\
Type gates applied & None & All at once & Per hop \\
\bottomrule
\end{tabular}
\end{table}

\paragraph{Trie size.}
The baseline and v1 build one trie from all paths ($O(d^L)$ entries).
v2 builds a new trie at each hop from only 1-hop neighbours ($O(d)$
entries per hop), but does so $L$ times.

\paragraph{Model passes.}
The dominant cost is the LLM forward pass ($\approx$6\,s on an A100 for
an 8B model).  Baseline and v1 each require one pass.  v2 requires up
to $L \times B$ passes (one per beam per hop), but in practice many
beams are pruned early.

\paragraph{When v2 wins.}
v2 is advantageous when:
\begin{itemize}
  \item The graph is large ($d^L$ is huge), so the baseline's trie
    is unwieldy.
  \item Type gates are aggressive, so filtering at each hop prunes
    more effectively than pre-filtering all paths.
  \item Beam search finds better paths than the baseline's single
    greedy/group-beam pass.
\end{itemize}

\paragraph{When baseline/v1 win.}
The baseline or v1 is advantageous when:
\begin{itemize}
  \item The graph is small ($d^L$ is manageable).
  \item Latency matters (one forward pass vs.\ many).
  \item Type gates are conservative (few paths removed), so the
    filtering overhead isn't justified.
\end{itemize}

%======================================================================
\section{Relationship Between the Algorithms}
%======================================================================

The three algorithms form a hierarchy:

\[
  \text{GCR Baseline} \;\supset\; \text{DCA-Trie v1} \;\subset\;
  \text{DCA-Trie v2}
\]

\begin{itemize}
  \item \textbf{v1 $\subset$ Baseline.}  v1's trie is a subset of the
    baseline's trie (filtered paths $\subseteq$ all paths).  v1
    generates from a smaller constraint set.
  \item \textbf{v2 is not a subset of either.}  v2's per-hop trie
    expansion explores a different set of paths than the baseline's
    full DFS.  Because v2 only enumerates 1-hop neighbours at each
    step (not all $L$-hop paths upfront), it can reach entities that
    the DFS might miss due to path explosion cutoffs.
  \item \textbf{All three use the same trie mechanism.}  The
    \texttt{prefix\_allowed\_tokens\_fn} callback is identical; only
    the trie's contents differ.
\end{itemize}

%======================================================================
\section{Discussion}
%======================================================================

\subsection{The Role of the TypeOracle}

The TypeOracle is \emph{not} a learned component.  It is a purely
symbolic rule engine that extracts type information from the KG's
ontology schema.  Its two gates enforce:

\begin{enumerate}
  \item \textbf{Local consistency} (range gate): the tail entity's type
    must be compatible with the relation's declared range.
  \item \textbf{Global consistency} (type gate): the terminal entity's
    type must be compatible with the expected answer types.
\end{enumerate}

The conservative default (admit when types are unknown) means the
oracle never over-filters when information is missing.  Its effectiveness
depends on the quality and coverage of the ontology schema.

\subsection{Limitations}

\begin{itemize}
  \item \textbf{Ontology coverage.}  The TypeOracle's hand-curated
    schema covers $\sim$38 Freebase relations.  Auto-mining from data
    triples extends coverage, but mined ranges may be noisy.
  \item \textbf{Evaluation metric.}  The standard Hits@1 metric uses
    substring matching ($\texttt{answer} \in \texttt{prediction}$),
    which is lenient.  Exact-match evaluation would give lower but more
    honest numbers.
  \item \textbf{v2 latency.}  Multiple forward passes make v2 slower
    than the single-pass baseline, despite smaller tries.
\end{itemize}

%======================================================================
\section*{Acknowledgements}
%======================================================================

This work builds on the GCR framework (Luo et al., 2025) and the DoG
decoding paradigm (Lin et al., 2025).

\end{document}
